\frfilename{mt2a6.tex}
\versiondate{10.11.14}
\copyrightdate{1994}

\def\chaptername{Lampiran}
\def\sectionname{Faktorisasi matriks}

\newsection{2A6}

Saya menggunakan beberapa halaman untuk membahas aljabar linear $\BbbR^r$
yang diperlukan bagi Bab 26. Saya hanya memberikan satu bukti, karena materi
ini dapat ditemukan dalam setiap buku ajar aljabar linear elementer; tetapi
saya rasa akan bermanfaat untuk meninjau gagasan-gagasan dasarnya dalam bahasa
yang saya gunakan untuk risalah ini.

\leader{2A6A}{Determinan}\cmmnt{ Kita perlu mengetahui hal-hal berikut
tentang determinan.}

\inset{(i) Setiap matriks real $r\times r$ $T$ mempunyai determinan real
$\det T$.}

\inset{(ii) Untuk setiap matriks $r\times r$ $S$ dan $T$,
$\det ST=\det S\det T$.}

\inset{(iii) Jika $T$ matriks diagonal, determinannya adalah hasil kali
entri-entri diagonalnya.}

\inset{(iv) Untuk setiap matriks $r\times r$ $T$, $\det T\trs=\det T$,
dengan $T\trs$ transpos dari $T$.}

\inset{(v) $\det T$ merupakan fungsi kontinu dari koefisien-koefisien
$T$.}

\cmmnt{\noindent Ada begitu banyak jalan untuk membahas topik ini sehingga
saya bahkan menghindari pemberian definisi `determinan'; saya mengajak Anda
memeriksa ingatan Anda, atau buku kesukaan Anda, untuk memastikan bahwa Anda
memang memahami fakta-fakta di atas.
}%end of comment

\leader{2A6B}{Keluarga ortonormal} Untuk $x=(\xi_1,\ldots,\xi_r)$,
$y=(\eta_1,\ldots,\eta_r)\in\BbbR^r$, tuliskan $x\dotproduct
y=\sum_{i=1}^r\xi_i\eta_i$; \cmmnt{tentu saja} $\|x\|$\cmmnt{, sebagaimana
didefinisikan dalam 1A2A,} adalah $\sqrt{x\dotproduct x}$.
\cmmnt{Ingat bahwa} $x_1,\ldots,x_k$ disebut {\bf ortonormal} jika
$x_i\dotproduct x_j=0$ untuk $i\ne j$, dan $1$ untuk $i=j$.
\cmmnt{Hasil-hasil yang kita perlukan di sini adalah:}

\inset{(i) Jika $x_1,\ldots,x_k$ merupakan vektor-vektor ortonormal dalam
$\BbbR^r$, dengan $k<r$, maka terdapat vektor $x_{k+1},\ldots,x_r$ dalam
$\BbbR^r$ sedemikian sehingga $x_1,\ldots,x_r$ ortonormal.}

\inset{(ii) Matriks $r\times r$ $P$ disebut {\bf ortogonal} jika $P\trs P$
adalah matriks identitas; secara ekuivalen, jika kolom-kolom $P$ ortonormal.}

\inset{(iii) Untuk matriks ortogonal $P$, $\det P$ harus bernilai
$\pm 1$\cmmnt{ (gabungkan (ii)--(iv) dari 2A6A)}.}
\inset{(iv) Jika $P$ ortogonal, maka $Px\dotproduct
Py\cmmnt{\mskip5mu =P\trs Px\dotproduct
y}=x\dotproduct y$ untuk semua $x$, $y\in\BbbR^r$.}

\inset{(v) Jika $P$ ortogonal, maka $P\trs=P^{-1}$ juga ortogonal.}

\inset{(vi) Jika $P$ dan $Q$ ortogonal, maka $PQ$ juga ortogonal.}

\leader{2A6C}{}\cmmnt{ Sekarang saya memberikan suatu proposisi yang tidak
selalu disertakan dalam penyajian elementer. Tentu saja ada banyak pendekatan
untuk hal ini; saya menawarkan pendekatan langsung.

\medskip

\noindent}{\bf Proposisi} Misalkan $T$ sebarang matriks real $r\times r$.
Maka $T$ dapat dinyatakan sebagai $PDQ$, dengan $P$ dan $Q$ matriks ortogonal
dan $D$ matriks diagonal berkoefisien tak negatif.

\proof{ Kita menggunakan induksi pada $r$.

\medskip

{\bf (a)} Jika $r=1$, maka $T=(\tau_{11})$. Tetapkan $D=(|\tau_{11}|)$,
$P=(1)$, dan $Q=(1)$ jika $\tau_{11}\ge 0$, serta $(-1)$ jika tidak.

\medskip

{\bf (b)(i)} Untuk langkah induksi ke $r+1\ge 2$, pertimbangkan bola satuan
$B=\{x:x\in\BbbR^{r+1},\,\|x\|\le 1\}$. Ini adalah himpunan tertutup dan
terbatas dalam $\BbbR^{r+1}$, sehingga kompak (2A2F). Pemetaan
$x\mapsto Tx:\BbbR^{r+1}\to\BbbR^{r+1}$ dan
$x\mapsto\|x\|:\BbbR^{r+1}\to\Bbb R$ kontinu, sehingga fungsi
$x\mapsto\|Tx\|:B\to\Bbb R$ terbatas dan mencapai batas-batasnya (2A2G),
dan terdapat $u\in B$ sedemikian sehingga $\|Tu\|\ge\|Tx\|$ untuk setiap
$x\in B$. Perhatikan bahwa $\|Tu\|$ harus merupakan norma $\|T\|$ dari $T$
sebagaimana didefinisikan dalam 262H. Tetapkan
$\delta=\|T\|=\|Tu\|$. Jika $\delta=0$, maka $T$ harus merupakan matriks
nol dan hasilnya trivial; jadi anggaplah $\delta>0$. Dalam kasus ini
$\|u\|$ harus tepat sama dengan $1$, sebab jika tidak, kita akan mempunyai
$u=\|u\|u'$ dengan $\|u'\|=1$ dan $\|Tu'\|>\|Tu\|$.

\medskip

\quad{\bf (ii)} Jika $x\in\BbbR^{r+1}$ dan $x\dotproduct u=0$, maka
$Tx\dotproduct Tu=0$. \Prf\Quer\ Jika tidak, tetapkan
$\gamma=Tx\dotproduct Tu\ne 0$. Pertimbangkan $y=u+\eta\gamma x$ untuk
$\eta>0$ yang kecil. Kita mempunyai

\Centerline{$\|y\|^2=y\dotproduct y
=u\dotproduct u+2\eta\gamma u\dotproduct x+\eta^2\gamma^2 x\dotproduct x
=\|u\|^2+\eta^2\gamma^2\|x\|^2
=1+\eta^2\gamma^2\|x\|^2$,}

\noindent sedangkan

\Centerline{$\|Ty\|^2=Ty\dotproduct Ty
=Tu\dotproduct Tu+2\eta\gamma Tu\dotproduct Tx
  +\eta^2\gamma^2 Tx\dotproduct Tx
=\delta^2+2\eta\gamma^2+\eta^2\gamma^2\|Tx\|^2$.}

\noindent Tetapi juga $\|Ty\|^2\le\delta^2\|y\|^2$ (2A4Fb), sehingga

\Centerline{$\delta^2+2\eta\gamma^2+\eta^2\gamma^2\|Tx\|^2
\le\delta^2(1+\eta^2\gamma^2\|x\|^2)$}

\noindent dan

\Centerline{$2\eta\gamma^2
\le\delta^2\eta^2\gamma^2\|x\|^2-\eta^2\gamma^2\|Tx\|^2$,}

\noindent yaitu

\Centerline{$2\le\eta(\delta^2\|x\|^2-\|Tx\|^2)$.}

\noindent Namun ini tentu tidak mungkin benar untuk semua $\eta>0$, jadi
kita memperoleh kontradiksi.\ \Bang\Qed

\medskip

\quad{\bf (iii)} Tetapkan $v=\delta^{-1}Tu$, sehingga $\|v\|=1$. Misalkan
$u_1,\ldots,u_{r+1}$ vektor-vektor ortonormal dengan $u_{r+1}=u$, dan
misalkan $Q_0$ matriks ortogonal $(r+1)\times(r+1)$ dengan kolom-kolom
$u_1,\ldots,u_{r+1}$; maka, dengan menuliskan $e_1,\ldots,e_{r+1}$ untuk
basis ortonormal baku dari $\BbbR^{r+1}$, kita mempunyai $Q_0e_i=u_i$ untuk
setiap $i$, dan $Q_0e_{r+1}=u$. Demikian pula, terdapat matriks ortogonal
$P_0$ sedemikian sehingga $P_0e_{r+1}=v$.

Tetapkan $T_1=P_0^{-1}TQ_0$. Maka

\Centerline{$T_1e_{r+1}=P_0^{-1}Tu=\delta P_0^{-1}v=\delta e_{r+1}$,}

\noindent sedangkan jika $x\dotproduct e_{r+1}=0$, maka
$Q_0x\dotproduct u=0$ (2A6B(iv)), sehingga

\Centerline{$T_1x\dotproduct e_{r+1}=P_0T_1x\dotproduct P_0e_{r+1}
=TQ_0x\dotproduct v=0$,}

\noindent menurut (ii). Ini berarti bahwa $T_1$ harus berbentuk

$$\Matrix{S&0\\0&\delta},$$

\noindent dengan $S$ matriks $r\times r$.

\medskip

\quad{\bf (iv)} Menurut hipotesis induksi, $S$ dapat dinyatakan sebagai
$\tilde P\tilde D\tilde Q$, dengan $\tilde P$ dan $\tilde Q$ matriks
ortogonal $r\times r$ dan $\tilde D$ matriks diagonal $r\times r$
berkoefisien tak negatif. Tetapkan

$$P_1=\Matrix{\tilde P&0\\0&1},\quad Q_1=\Matrix{\tilde Q&0\\0&1},\quad
D=\Matrix{\tilde D&0\\0&\delta}.$$

\noindent Maka $P_1$ dan $Q_1$ ortogonal dan $D$ diagonal berkoefisien tak
negatif, serta $P_1DQ_1=T_1$. Sekarang tetapkan

\Centerline{$P=P_0P_1$,\quad $Q=Q_1Q_0^{-1}$,}
\noindent sehingga $P$ dan $Q$ ortogonal (2A6B(v)--(vi)) dan

\Centerline{$PDQ=P_0P_1DQ_1Q_0^{-1}=P_0T_1Q_0^{-1}=T$.}

\noindent Jadi langkah induksi selesai.
}%end of proof of 2A6C

\frnewpage

